\begin{tikzpicture}[
    font=\small,
    box/.style={
        draw,
        rounded corners=2pt,
        align=center,
        minimum height=9mm,
        text width=30mm,
        inner sep=3pt
    },
    input/.style={
        box,
        fill=blue!8
    },
    rgrid/.style={
        box,
        fill=cyan!10
    },
    behav/.style={
        box,
        fill=orange!12
    },
    cash/.style={
        box,
        fill=green!10
    },
    price/.style={
        box,
        fill=purple!8
    },
    dec/.style={
        box,
        fill=gray!8
    },
    arr/.style={
        ->,
        thick
    },
    darr/.style={
        ->,
        thick,
        dashed
    },
    lab/.style={
        font=\scriptsize
    },
    panel/.style={
        draw=gray!40,
        rounded corners=4pt,
        inner sep=8pt,
        fill=gray!4
    },
    ptitle/.style={
        anchor=south west,
        font=\scriptsize\scshape
    }
]

% ==================================================
% 1. Input og rentemodel
% ==================================================

\node[input] (kurve) at (0,0)
    {Dagens\\rentekurve};

\node[input] (tsm) at (5.8,0)
    {Rentestruktur-\\model};

\node[rgrid] (grid) at (11.6,0)
    {Rentegitter eller\\simulerede rentestier};

\draw[arr] (kurve) -- (tsm);
\draw[arr] (tsm) -- (grid);


% ==================================================
% 2. Låntageradfærd og cash flows
% ==================================================

\node[behav] (lys) at (0,-4.0) {%
    Refinansierings-\\kurve\\
    {\footnotesize
    \(z^{\text{refi}}_{\tau}
      =
      z^{\text{swap}}_{\tau}
      +
      LYS_{\tau}\)}
};

\node[behav] (gain) at (5.8,-4.0) {%
    Konverterings-\\gevinst\\
    \(Gain_t\)
};

\node[behav] (cpr) at (11.6,-4.0) {%
    \(CPR_t\) og \(UDT_t\)
};

\node[behav, text width=35mm] (burn) at (5.8,-7.0) {%
    Burnout \(B_t\)\\
    {\scriptsize den tilbageværende låntagerpulje}
};

\node[cash] (cf) at (11.6,-7.0) {%
    Cash flows\\
    på hver sti
};

% Renten påvirker låntagernes konverteringsincitament
\draw[
    arr,
    rounded corners=6pt
]
    (grid.south)
    --
    (11.6,-2.1)
    --
    node[
        lab,
        fill=white,
        inner sep=1.5pt
    ] {%
        \textbf{kanal 2:}
        renten påvirker konverteringsincitamentet
    }
    (5.8,-2.1)
    --
    (gain.north);

% Refinansieringskurven bestemmer sammen med renten Gain
\draw[arr] (lys) -- (gain);

% Gain og den eksisterende pulje bestemmer CPR
\draw[arr] (gain) -- (cpr);

\draw[
    darr,
    rounded corners=4pt
]
    (burn.north)
    -- (5.8,-5.45)
    -- node[
        lab,
        fill=white,
        inner sep=1.5pt,
        above=1pt
    ] {%
        påvirker næste termins CPR
    }
    (11.6,-5.45)
    -- (cpr.south);

% CPR bestemmer cash flows
\draw[arr] (cpr) -- (cf);

% CPR opdaterer burnout
\draw[
    arr
]
    (cpr.south west)
    to[out=245,in=10]
    node[
        lab,
        fill=white,
        inner sep=1.5pt,
        pos=0.55,
        below=1pt
    ] {%
        burnout opdateres
    }
    (burn.east);


% ==================================================
% 3. Diskontering og modelpris
% ==================================================

\node[price] (mp) at (5.8,-10.0) {%
    Modelpris\\
    \(P_{\text{model}}(OAS)\)
};

\node[price] (disc) at (11.6,-10.0) {%
    Diskontering pr.\ sti\\
    {\footnotesize
    \(z^{\text{swap}}_{\tau}+OAS\)}
};

% Cash flows diskonteres
\draw[arr] (cf) -- (disc);

% De diskonterede cash flows giver modelprisen
\draw[arr] (disc) -- (mp);

% Renten påvirker også diskonteringen direkte
\draw[
    arr,
    rounded corners=6pt
]
    (grid.east)
    --
    (15.4,0)
    --
    (15.4,-10.0)
    --
    (disc.east);

\node[
    lab,
    rotate=90,
    anchor=south
] at (15.4,-5.0) {%
    \textbf{kanal 1:}
    renten påvirker diskonteringen \(DF_t^{(m)}\)
};


% ==================================================
% 4. Kalibrering til markedskurs
% ==================================================

\node[price] (mk) at (0,-12.8) {%
    Markedskurs\\
    \(P_{\text{market}}\)
};

\node[dec] (cmp) at (5.8,-12.8) {%
    Er priserne ens?\\
    {\footnotesize
    \(P_{\text{model}}
      \stackrel{?}{=}
      P_{\text{market}}\)}
};

\draw[arr] (mp) -- (cmp);
\draw[arr] (mk) -- (cmp);

% Hvis priserne ikke er ens, justeres OAS
\draw[
    darr,
    rounded corners=6pt
]
    (cmp.east)
    --
    node[
        lab,
        fill=white,
        inner sep=1.5pt,
        pos=0.68,
        above=2pt,
        align=center
    ] {%
        \textbf{Nej:}
        justér OAS og kør modellen igen
    }
    (11.6,-12.8)
    --
    (disc.south);


% ==================================================
% 5. Resultat og nøgletal
% ==================================================

\node[price] (ok) at (5.8,-15.4) {%
    \textbf{Ja: OAS fundet}\\
    {\scriptsize kalibreret model}
};

\node[
    price,
    text width=44mm
] (ngl) at (12.0,-15.4) {%
    Nøgletal: varighed, konveksitet,\\
    CPR-, LYS- og OAS-risk
};

\draw[arr] (cmp) -- (ok);
\draw[arr] (ok) -- (ngl);


% ==================================================
% Paneler
% ==================================================

\begin{scope}[on background layer]

\node[
    panel,
    fit=(kurve)(tsm)(grid)
] (pA) {};

\node[
    panel,
    fit=(lys)(gain)(cpr)(burn)(cf)
] (pB) {};

\node[
    panel,
    fit=(mp)(disc)
] (pC) {};

\node[
    panel,
    fit=(mk)(cmp)
] (pD) {};

\node[
    panel,
    fit=(ok)(ngl)
] (pE) {};

\end{scope}

\node[ptitle] at ([yshift=2pt]pA.north west)
    {1.\ Input og rentemodel};

\node[ptitle] at ([yshift=2pt]pB.north west)
    {2.\ Låntageradfærd og cash flows};

\node[ptitle] at ([yshift=2pt]pC.north west)
    {3.\ Diskontering og modelpris};

\node[ptitle] at ([yshift=2pt]pD.north west)
    {4.\ Kalibrering til markedskurs};

\node[ptitle, anchor=south east] at ([yshift=2pt]pE.north east)
    {5.\ Nøgletal};

\end{tikzpicture}%
